\documentclass[12pt,a4paper]{article}

\usepackage[T2A]{fontenc}
\usepackage[utf8]{inputenc}
\usepackage[russian]{babel}
\usepackage{amsmath,amssymb,amsthm}
\usepackage{enumitem}
\usepackage{array}
\usepackage{booktabs}
\usepackage{mathtools}
\usepackage{geometry}
\usepackage{hyperref}

\geometry{left=2.2cm,right=2.2cm,top=2cm,bottom=2cm}

\newtheorem{definition}{Определение}
\newtheorem{theorem}{Теорема}
\newtheorem{remark}{Замечание}
\newtheorem{example}{Пример}

\title{Билет №25 \\ \small Многопроцессорные и многомашинные комплексы. Вычислительные кластеры. Распределенные системы. (ИТМО). \\ Назначение, архитектура и принципы построения информационно-вычислительных сетей (ИВС). Локальные и глобальные ИВС, технические и программные средства объединения различных сетей. (СпБГУ)}
\author{Сериков Владислав Александрович}
\date{4 мая 2026}

\begin{document}

\maketitle
\tableofcontents
\newpage

\section{Базовые понятия}

\begin{definition}
\textbf{Вычислительная система} --- совокупность аппаратных и программных средств, предназначенных для хранения, передачи, обработки данных и управления вычислительными процессами.
\end{definition}

\begin{definition}
\textbf{Параллельная вычислительная система} --- система, в которой несколько вычислительных устройств одновременно выполняют части одной или нескольких задач.
\end{definition}

\begin{definition}
\textbf{Распределённая система} --- совокупность независимых компьютеров, связанных сетью и воспринимаемых пользователем или прикладной программой как единая согласованная система.
\end{definition}

Ключевые свойства распределённых и параллельных систем:
\begin{itemize}
    \item \textbf{параллелизм} --- возможность одновременного выполнения вычислений;
    \item \textbf{масштабируемость} --- способность сохранять эффективность при росте числа узлов, пользователей или объёма данных;
    \item \textbf{отказоустойчивость} --- способность продолжать работу при отказе части компонентов;
    \item \textbf{прозрачность} --- сокрытие от пользователя факта распределённости ресурсов;
    \item \textbf{совместное использование ресурсов} --- доступ к общим данным, устройствам, сервисам;
    \item \textbf{гетерогенность} --- возможность объединения различных аппаратных платформ, ОС, сетей и протоколов.
\end{itemize}

\section{Многопроцессорные вычислительные системы}

\subsection{Понятие многопроцессорной системы}

\begin{definition}
\textbf{Многопроцессорная система} --- вычислительная система, содержащая несколько процессоров, работающих под управлением общей операционной системы и, как правило, имеющих доступ к общей памяти.
\end{definition}

Основная идея многопроцессорной системы состоит в том, что одна задача может быть разделена на несколько параллельно выполняемых частей, либо несколько независимых задач могут выполняться одновременно.

\subsection{Классификация по Флинну}

Классическая классификация вычислительных архитектур основана на числе потоков команд и потоков данных.

\begin{center}
\begin{tabular}{lll}
\toprule
Тип & Расшифровка & Характеристика \\
\midrule
SISD & Single Instruction Single Data & Один поток команд, один поток данных \\
SIMD & Single Instruction Multiple Data & Один поток команд, много потоков данных \\
MISD & Multiple Instruction Single Data & Много потоков команд, один поток данных \\
MIMD & Multiple Instruction Multiple Data & Много потоков команд, много потоков данных \\
\bottomrule
\end{tabular}
\end{center}

\textbf{SISD} соответствует классическому однопроцессорному компьютеру.

\textbf{SIMD} используется в векторных процессорах, GPU, мультимедийных расширениях процессоров. Один управляющий поток выполняет одну инструкцию над множеством элементов данных.

\textbf{MISD} практически редко встречается в чистом виде, но идея используется в отказоустойчивых системах, где один поток данных обрабатывается несколькими независимыми устройствами.

\textbf{MIMD} --- наиболее распространённая модель для многопроцессорных систем, кластеров и распределённых систем.

\subsection{SMP-системы}

\begin{definition}
\textbf{SMP} --- symmetric multiprocessing, симметричная многопроцессорная система, в которой несколько однотипных процессоров имеют равноправный доступ к общей памяти и устройствам ввода-вывода.
\end{definition}

Схематически SMP можно представить так:
\[
\text{CPU}_1,\text{CPU}_2,\ldots,\text{CPU}_n
\longleftrightarrow
\text{общая шина или коммутатор}
\longleftrightarrow
\text{общая память}.
\]

Преимущества SMP:
\begin{itemize}
    \item простая модель программирования за счёт общей памяти;
    \item единый образ операционной системы;
    \item удобное взаимодействие потоков через общие переменные;
    \item хорошая производительность при умеренном числе процессоров.
\end{itemize}

Недостатки:
\begin{itemize}
    \item конкуренция за доступ к памяти;
    \item ограниченная масштабируемость;
    \item необходимость обеспечения когерентности кэшей;
    \item рост сложности аппаратной синхронизации.
\end{itemize}

\subsection{NUMA-системы}

\begin{definition}
\textbf{NUMA} --- non-uniform memory access, архитектура с неоднородным доступом к памяти, в которой память физически распределена между узлами, но логически образует единое адресное пространство.
\end{definition}

В NUMA-системе каждый процессор быстрее обращается к своей локальной памяти и медленнее --- к памяти других узлов:
\[
T_{\text{local}} < T_{\text{remote}}.
\]

Преимущества NUMA:
\begin{itemize}
    \item лучшая масштабируемость по сравнению с SMP;
    \item сохранение общей адресной модели;
    \item возможность построения систем с большим числом процессоров.
\end{itemize}

Недостатки:
\begin{itemize}
    \item производительность зависит от размещения данных;
    \item требуется NUMA-aware планировщик ОС;
    \item сложнее программировать высокоэффективные приложения.
\end{itemize}

\subsection{Когерентность кэшей}

\begin{definition}
\textbf{Когерентность кэшей} --- свойство многопроцессорной системы, гарантирующее, что разные процессоры видят согласованные значения одной и той же ячейки памяти.
\end{definition}

Проблема возникает из-за того, что каждый процессор имеет собственный кэш. Если два процессора кэшируют одну переменную $x$, а один из них изменяет её, другой может продолжать видеть устаревшее значение.

Основные подходы:
\begin{itemize}
    \item \textbf{write-through} --- запись сразу в кэш и основную память;
    \item \textbf{write-back} --- запись сначала в кэш, затем в память при вытеснении;
    \item \textbf{snooping} --- процессоры отслеживают операции на общей шине;
    \item \textbf{directory-based coherence} --- информация о копиях блоков хранится в специальных каталогах.
\end{itemize}

Типичный протокол когерентности --- MESI. Блок кэша может находиться в одном из состояний:
\begin{itemize}
    \item $M$ --- modified, блок изменён и отличается от памяти;
    \item $E$ --- exclusive, блок совпадает с памятью и есть только в одном кэше;
    \item $S$ --- shared, блок может быть в нескольких кэшах;
    \item $I$ --- invalid, блок недействителен.
\end{itemize}

\subsection{Синхронизация в многопроцессорных системах}

При параллельном выполнении потоков возникает необходимость синхронизации доступа к общим данным.

Основные средства:
\begin{itemize}
    \item мьютексы;
    \item семафоры;
    \item мониторы;
    \item барьеры;
    \item атомарные операции;
    \item спин-блокировки.
\end{itemize}

\begin{definition}
\textbf{Критическая секция} --- участок программы, в котором выполняется доступ к разделяемому ресурсу и который не должен одновременно выполняться несколькими потоками.
\end{definition}

Требования к решению задачи критической секции:
\begin{enumerate}
    \item \textbf{взаимное исключение}: не более одного процесса в критической секции;
    \item \textbf{прогресс}: если критическая секция свободна, выбор следующего процесса не откладывается бесконечно;
    \item \textbf{ограниченное ожидание}: процесс не должен ждать входа в критическую секцию бесконечно долго.
\end{enumerate}

\section{Многомашинные вычислительные комплексы}

\subsection{Понятие многомашинного комплекса}

\begin{definition}
\textbf{Многомашинный вычислительный комплекс} --- совокупность нескольких самостоятельных компьютеров, соединённых каналами связи и совместно решающих вычислительные, информационные или управляющие задачи.
\end{definition}

В отличие от многопроцессорной системы, многомашинный комплекс обычно состоит из машин, каждая из которых имеет:
\begin{itemize}
    \item собственную оперативную память;
    \item собственную операционную систему;
    \item собственные устройства ввода-вывода;
    \item собственный сетевой интерфейс.
\end{itemize}

Обмен данными между машинами выполняется по сети с помощью сообщений.

\subsection{Сравнение многопроцессорных и многомашинных систем}

\begin{center}
\begin{tabular}{p{4cm}p{5cm}p{5cm}}
\toprule
Критерий & Многопроцессорная система & Многомашинный комплекс \\
\midrule
Память & Общая или логически общая & Раздельная \\
ОС & Обычно единая & На каждом узле своя ОС \\
Связь & Через память, шину, interconnect & Через сеть сообщений \\
Масштабируемость & Ограничена памятью и interconnect & Выше, особенно в кластерах \\
Программирование & Потоки, общая память & Сообщения, RPC, MPI \\
Отказы & Отказ узла часто критичен & Возможна изоляция отказа \\
\bottomrule
\end{tabular}
\end{center}

\subsection{Модели взаимодействия}

Для многомашинных систем характерны две основные модели взаимодействия:

\begin{enumerate}
    \item \textbf{Передача сообщений}: процессы обмениваются сообщениями через сеть.
    \[
    P_i \xrightarrow{\text{message}} P_j.
    \]

    \item \textbf{Распределённая общая память}: программная система имитирует общее адресное пространство поверх физически распределённой памяти.
\end{enumerate}

Передача сообщений более явно отражает устройство системы, но требует от программиста управления обменами. Распределённая общая память удобнее, но сложнее реализуется эффективно.

\section{Вычислительные кластеры}

\subsection{Определение и назначение}

\begin{definition}
\textbf{Вычислительный кластер} --- группа соединённых сетью компьютеров, работающих совместно и воспринимаемых как единый вычислительный ресурс.
\end{definition}

Кластеры применяются для:
\begin{itemize}
    \item высокопроизводительных вычислений;
    \item обработки больших данных;
    \item повышения доступности сервисов;
    \item балансировки нагрузки;
    \item научного моделирования;
    \item машинного обучения;
    \item веб-сервисов и облачных платформ.
\end{itemize}

\subsection{Типы кластеров}

\subsubsection{Высокопроизводительные кластеры}

\begin{definition}
\textbf{HPC-кластер} --- кластер, предназначенный для выполнения ресурсоёмких вычислительных задач с минимальным временем решения.
\end{definition}

Примеры задач:
\begin{itemize}
    \item численное моделирование;
    \item расчёт матриц;
    \item аэродинамика;
    \item квантовая химия;
    \item биоинформатика.
\end{itemize}

\subsubsection{Кластеры высокой доступности}

\begin{definition}
\textbf{HA-кластер} --- high availability cluster, кластер, предназначенный для обеспечения непрерывности работы сервиса при отказе отдельных узлов.
\end{definition}

Обычно используется резервирование:
\[
\text{основной узел} + \text{резервный узел}.
\]

Если основной узел отказывает, резервный принимает его функции. Такой процесс называется \textbf{failover}.

\subsubsection{Кластеры балансировки нагрузки}

\begin{definition}
\textbf{Кластер балансировки нагрузки} --- система, распределяющая входящие запросы между несколькими серверами.
\end{definition}

Пример:
\[
\text{клиенты}
\rightarrow
\text{балансировщик}
\rightarrow
\{\text{сервер}_1,\text{сервер}_2,\ldots,\text{сервер}_n\}.
\]

\subsection{Архитектура вычислительного кластера}

Типичный кластер включает:
\begin{itemize}
    \item вычислительные узлы;
    \item управляющий узел;
    \item сеть управления;
    \item высокоскоростную сеть обмена данными;
    \item систему хранения;
    \item планировщик заданий;
    \item программную среду параллельного выполнения.
\end{itemize}

\subsection{Планирование заданий в кластере}

Планировщик принимает задания пользователей и распределяет их по узлам.

Основные функции:
\begin{itemize}
    \item постановка заданий в очередь;
    \item выделение ресурсов;
    \item запуск задач;
    \item контроль выполнения;
    \item сбор статуса;
    \item завершение или перезапуск задач при ошибках.
\end{itemize}

Критерии планирования:
\begin{itemize}
    \item минимизация времени ожидания;
    \item максимизация загрузки ресурсов;
    \item справедливое распределение ресурсов;
    \item соблюдение приоритетов;
    \item учёт ограничений по памяти, CPU, GPU и сети.
\end{itemize}

\subsection{MPI как модель программирования кластеров}

\begin{definition}
\textbf{MPI} --- Message Passing Interface, стандарт программного интерфейса для обмена сообщениями между процессами в параллельных программах.
\end{definition}

Основные операции MPI:
\begin{itemize}
    \item \texttt{MPI\_Init} --- инициализация;
    \item \texttt{MPI\_Comm\_size} --- получение числа процессов;
    \item \texttt{MPI\_Comm\_rank} --- получение номера процесса;
    \item \texttt{MPI\_Send} --- отправка сообщения;
    \item \texttt{MPI\_Recv} --- получение сообщения;
    \item \texttt{MPI\_Bcast} --- широковещательная рассылка;
    \item \texttt{MPI\_Reduce} --- свёртка данных;
    \item \texttt{MPI\_Finalize} --- завершение.
\end{itemize}

\begin{example}
Минимальный пример распределения суммы массива между четырьмя процессами:
\[
A = [1,2,3,4,5,6,7,8].
\]
Пусть есть $4$ процесса. Каждый получает по два элемента:
\[
P_0: [1,2],\quad P_1: [3,4],\quad P_2: [5,6],\quad P_3: [7,8].
\]
Локальные суммы:
\[
s_0=3,\quad s_1=7,\quad s_2=11,\quad s_3=15.
\]
Глобальная сумма:
\[
S = s_0+s_1+s_2+s_3=36.
\]
\end{example}

\section{Оценка ускорения параллельных вычислений}

\subsection{Ускорение и эффективность}

\begin{definition}
\textbf{Ускорение} параллельной программы на $p$ процессорах:
\[
S_p = \frac{T_1}{T_p},
\]
где $T_1$ --- время работы на одном процессоре, $T_p$ --- время работы на $p$ процессорах.
\end{definition}

\begin{definition}
\textbf{Эффективность}:
\[
E_p = \frac{S_p}{p}.
\]
\end{definition}

Если $E_p=1$, то ресурсы используются идеально. На практике:
\[
0 < E_p < 1.
\]

\subsection{Закон Амдала}

\begin{theorem}[Закон Амдала]
Пусть доля последовательной части программы равна $\alpha$, где $0 \leq \alpha \leq 1$, а оставшаяся доля $1-\alpha$ идеально распараллеливается на $p$ процессорах. Тогда максимальное ускорение:
\[
S_p = \frac{1}{\alpha + \frac{1-\alpha}{p}}.
\]
При $p \to \infty$ предельное ускорение равно:
\[
S_{\infty} = \frac{1}{\alpha}.
\]
\end{theorem}

\begin{proof}
Пусть время выполнения программы на одном процессоре равно $T_1$.

Последовательная часть занимает:
\[
T_{\text{seq}} = \alpha T_1.
\]

Параллельная часть занимает:
\[
T_{\text{par}} = (1-\alpha)T_1.
\]

Если параллельная часть идеально делится между $p$ процессорами, то её время выполнения становится:
\[
\frac{(1-\alpha)T_1}{p}.
\]

Итого время выполнения на $p$ процессорах:
\[
T_p = \alpha T_1 + \frac{(1-\alpha)T_1}{p}
= T_1\left(\alpha + \frac{1-\alpha}{p}\right).
\]

Ускорение:
\[
S_p = \frac{T_1}{T_p}
= \frac{T_1}{T_1\left(\alpha + \frac{1-\alpha}{p}\right)}
= \frac{1}{\alpha + \frac{1-\alpha}{p}}.
\]

Переходя к пределу при $p \to \infty$, получаем:
\[
\lim_{p\to\infty} S_p
=
\lim_{p\to\infty}
\frac{1}{\alpha + \frac{1-\alpha}{p}}
=
\frac{1}{\alpha}.
\]

Теорема доказана.
\end{proof}

\begin{example}
Пусть $10\%$ программы не распараллеливается, то есть $\alpha=0.1$.
На $p=8$ процессорах:
\[
S_8 = \frac{1}{0.1+\frac{0.9}{8}}
= \frac{1}{0.1+0.1125}
= \frac{1}{0.2125}
\approx 4.71.
\]
Даже при бесконечном числе процессоров:
\[
S_{\infty}=\frac{1}{0.1}=10.
\]
\end{example}

\subsection{Закон Густафсона}

Закон Амдала рассматривает фиксированный размер задачи. Закон Густафсона отражает ситуацию, когда с ростом числа процессоров увеличивают размер задачи.

Если последовательная доля равна $\alpha$, то масштабируемое ускорение:
\[
S_p = p - \alpha(p-1).
\]

Идея: большее число процессоров позволяет решать задачу большего размера за то же время.

\section{Распределённые системы}

\subsection{Основные свойства}

Распределённая система состоит из узлов, взаимодействующих по сети. Узлы могут отказывать независимо, сообщения могут задерживаться, теряться или приходить в другом порядке.

Ключевые свойства:
\begin{itemize}
    \item отсутствие глобальных часов;
    \item независимые отказы компонентов;
    \item асинхронность взаимодействия;
    \item конкуренция за общие ресурсы;
    \item необходимость согласования состояния;
    \item проблемы репликации и согласованности данных.
\end{itemize}

\subsection{Прозрачность распределённых систем}

Виды прозрачности:
\begin{itemize}
    \item \textbf{прозрачность доступа}: одинаковый доступ к локальным и удалённым ресурсам;
    \item \textbf{прозрачность местоположения}: пользователь не знает, где расположен ресурс;
    \item \textbf{прозрачность миграции}: ресурс может перемещаться без изменения интерфейса;
    \item \textbf{прозрачность репликации}: несколько копий ресурса выглядят как один ресурс;
    \item \textbf{прозрачность отказов}: система скрывает отказы компонентов;
    \item \textbf{прозрачность масштабирования}: расширение системы не требует изменения приложений.
\end{itemize}

\subsection{Архитектурные модели}

\subsubsection{Клиент-серверная модель}

Клиент отправляет запрос, сервер возвращает ответ:
\[
\text{client} \rightarrow \text{request} \rightarrow \text{server},
\]
\[
\text{client} \leftarrow \text{response} \leftarrow \text{server}.
\]

Преимущества:
\begin{itemize}
    \item простота;
    \item централизованное управление;
    \item удобная защита данных.
\end{itemize}

Недостатки:
\begin{itemize}
    \item сервер может стать узким местом;
    \item отказ сервера нарушает работу системы;
    \item ограниченная масштабируемость без репликации.
\end{itemize}

\subsubsection{Многоуровневая архитектура}

Типичный вариант:
\[
\text{клиент} \rightarrow \text{сервер приложений} \rightarrow \text{сервер БД}.
\]

Преимущества:
\begin{itemize}
    \item разделение ответственности;
    \item масштабирование отдельных уровней;
    \item повышение сопровождаемости.
\end{itemize}

\subsubsection{Peer-to-peer}

\begin{definition}
\textbf{P2P-система} --- распределённая система, в которой узлы одновременно могут выступать и клиентами, и серверами.
\end{definition}

Преимущества:
\begin{itemize}
    \item отсутствие единой центральной точки отказа;
    \item естественная масштабируемость;
    \item распределение нагрузки между участниками.
\end{itemize}

Недостатки:
\begin{itemize}
    \item сложность поиска ресурсов;
    \item трудность управления безопасностью;
    \item нестабильность состава узлов.
\end{itemize}

\subsection{Модели времени и событий}

В распределённой системе невозможно в общем случае полагаться на идеально синхронизированные физические часы. Поэтому вводятся логические часы.

\begin{definition}
Отношение \textbf{happened-before}, обозначаемое $\rightarrow$, определяется так:
\begin{enumerate}
    \item если события $a$ и $b$ происходят в одном процессе и $a$ раньше $b$, то $a \rightarrow b$;
    \item если $a$ --- отправка сообщения, а $b$ --- получение этого сообщения, то $a \rightarrow b$;
    \item если $a \rightarrow b$ и $b \rightarrow c$, то $a \rightarrow c$.
\end{enumerate}
\end{definition}

Если ни $a \rightarrow b$, ни $b \rightarrow a$, то события называются конкурентными.

\subsection{Алгоритм логических часов Лэмпорта}

\subsubsection{Идея}

Каждый процесс $P_i$ хранит счётчик $C_i$. Счётчик увеличивается при каждом локальном событии и передаётся вместе с сообщениями.

\subsubsection{Шаги алгоритма}

\begin{enumerate}
    \item Каждый процесс $P_i$ инициализирует свой счётчик:
    \[
    C_i := 0.
    \]

    \item Перед каждым локальным событием процесс увеличивает счётчик:
    \[
    C_i := C_i + 1.
    \]

    \item При отправке сообщения $m$ процесс прикрепляет к нему метку времени:
    \[
    timestamp(m) := C_i.
    \]

    \item При получении сообщения с меткой $t$ процесс обновляет счётчик:
    \[
    C_j := \max(C_j,t)+1.
    \]
\end{enumerate}

\subsubsection{Минимальный пример}

Пусть есть два процесса $P_1$ и $P_2$.

\[
P_1: C_1=0,\qquad P_2: C_2=0.
\]

Процесс $P_1$ выполняет локальное событие:
\[
C_1=1.
\]

Затем $P_1$ отправляет сообщение процессу $P_2$:
\[
C_1=2,\quad timestamp=2.
\]

Процесс $P_2$ до получения сообщения имеет:
\[
C_2=0.
\]

После получения:
\[
C_2=\max(0,2)+1=3.
\]

Следовательно, получение сообщения получает логическое время $3$, что больше времени отправки $2$.

\subsection{Выбор координатора: алгоритм Bully}

\subsubsection{Постановка задачи}

В распределённой системе часто нужен координатор: например, для распределения ресурсов, управления блокировками или организации репликации.

Пусть процессы имеют уникальные идентификаторы. Чем больше идентификатор, тем выше приоритет.

\subsubsection{Шаги алгоритма Bully}

\begin{enumerate}
    \item Процесс $P_i$ обнаруживает, что координатор не отвечает.
    \item $P_i$ отправляет сообщение \texttt{ELECTION} всем процессам с идентификаторами больше, чем у него.
    \item Если никто не отвечает, $P_i$ объявляет себя координатором и рассылает сообщение \texttt{COORDINATOR}.
    \item Если хотя бы один процесс с большим идентификатором ответил \texttt{OK}, то $P_i$ прекращает свою попытку выбора и ждёт результата.
    \item Процесс с большим идентификатором, получив \texttt{ELECTION}, сам начинает выборы среди процессов с ещё большими идентификаторами.
    \item В итоге координатором становится работоспособный процесс с максимальным идентификатором.
\end{enumerate}

\subsubsection{Минимальный пример}

Пусть есть процессы:
\[
P_1,\ P_2,\ P_3,\ P_4,\ P_5.
\]
Идентификатор равен номеру процесса. Координатором был $P_5$, но он отказал.

Процесс $P_2$ обнаружил отказ:
\[
P_2 \rightarrow P_3,P_4,P_5: \texttt{ELECTION}.
\]

$P_3$ и $P_4$ отвечают:
\[
P_3,P_4 \rightarrow P_2: \texttt{OK}.
\]

$P_4$ начинает свои выборы:
\[
P_4 \rightarrow P_5: \texttt{ELECTION}.
\]

$P_5$ не отвечает. Тогда:
\[
P_4 \rightarrow P_1,P_2,P_3: \texttt{COORDINATOR}.
\]

Итог: новым координатором становится $P_4$.

\section{Согласованность, доступность и разделение сети}

\subsection{Репликация данных}

\begin{definition}
\textbf{Репликация} --- хранение нескольких копий одного и того же объекта данных на разных узлах.
\end{definition}

Цели репликации:
\begin{itemize}
    \item повышение доступности;
    \item снижение задержек чтения;
    \item распределение нагрузки;
    \item отказоустойчивость.
\end{itemize}

Проблема: при изменении одной копии нужно согласовать остальные.

\subsection{Модели согласованности}

\begin{itemize}
    \item \textbf{Строгая согласованность}: чтение всегда возвращает результат самой последней записи.
    \item \textbf{Последовательная согласованность}: результат эквивалентен некоторому последовательному порядку операций всех процессов.
    \item \textbf{Линеаризуемость}: каждая операция выглядит выполненной мгновенно в некоторой точке между её вызовом и завершением.
    \item \textbf{Eventual consistency}: если обновления прекратятся, все реплики в конечном итоге придут к одному состоянию.
\end{itemize}

\subsection{CAP-теорема}

\begin{definition}
В контексте распределённых хранилищ обычно выделяют три свойства:
\begin{itemize}
    \item \textbf{Consistency} --- все клиенты видят согласованные данные;
    \item \textbf{Availability} --- каждый запрос к исправному узлу завершается ответом;
    \item \textbf{Partition tolerance} --- система продолжает работу при разделении сети.
\end{itemize}
\end{definition}

\begin{theorem}[CAP-теорема, неформальная формулировка]
В асинхронной распределённой системе хранения данных при возможном разделении сети невозможно одновременно гарантировать согласованность, доступность и устойчивость к разделению сети.
\end{theorem}

\begin{proof}
Рассмотрим систему из двух узлов $A$ и $B$, на которых хранится реплика одного объекта $x$.

Пусть начальное значение:
\[
x=0.
\]

Предположим, произошло разделение сети: узлы $A$ и $B$ не могут обмениваться сообщениями.

Клиент $C_1$ обращается к узлу $A$ и записывает:
\[
x:=1.
\]

Если система требует доступности, узел $A$ должен завершить операцию записи успешно, не ожидая связи с $B$.

Теперь клиент $C_2$ обращается к узлу $B$ и читает $x$.

Так как сеть разделена, узел $B$ не знает о записи $x:=1$ на узле $A$.

Есть два варианта:
\begin{enumerate}
    \item Узел $B$ отвечает значением $0$. Тогда доступность сохранена, но согласованность нарушена.
    \item Узел $B$ отказывается отвечать до восстановления связи. Тогда согласованность можно сохранить, но доступность нарушена.
\end{enumerate}

Следовательно, при наличии разделения сети нельзя одновременно гарантировать и доступность, и согласованность. Поскольку устойчивость к разделению предполагает, что система должна рассматривать такой случай, все три свойства одновременно недостижимы.

Теорема доказана в неформальной модели.
\end{proof}

\section{Информационно-вычислительные сети}

\subsection{Назначение ИВС}

\begin{definition}
\textbf{Информационно-вычислительная сеть} --- совокупность компьютеров, серверов, коммуникационного оборудования, каналов связи и программных средств, обеспечивающих совместную обработку, хранение и передачу информации.
\end{definition}

Назначение ИВС:
\begin{itemize}
    \item обмен данными между пользователями и программами;
    \item совместное использование вычислительных ресурсов;
    \item доступ к базам данных и файловым хранилищам;
    \item организация корпоративных информационных систем;
    \item удалённое управление;
    \item обеспечение доступа к Интернету;
    \item поддержка распределённых приложений.
\end{itemize}

\subsection{Основные компоненты ИВС}

\begin{enumerate}
    \item \textbf{Абонентские системы}: рабочие станции, серверы, мобильные устройства.
    \item \textbf{Коммуникационные узлы}: коммутаторы, маршрутизаторы, точки доступа.
    \item \textbf{Каналы связи}: медные кабели, оптоволокно, радиоканалы.
    \item \textbf{Сетевые интерфейсы}: сетевые карты, Wi-Fi-адаптеры.
    \item \textbf{Сетевое ПО}: драйверы, протоколы, службы, сетевые ОС.
    \item \textbf{Прикладные сервисы}: DNS, DHCP, web, email, файловые службы.
\end{enumerate}

\subsection{Требования к ИВС}

\begin{itemize}
    \item производительность;
    \item надёжность;
    \item масштабируемость;
    \item безопасность;
    \item управляемость;
    \item совместимость;
    \item отказоустойчивость;
    \item экономическая эффективность.
\end{itemize}

\section{Архитектура ИВС}

\subsection{Уровневая организация}

Сетевые системы строятся по уровневому принципу. Каждый уровень:
\begin{itemize}
    \item предоставляет сервис верхнему уровню;
    \item использует сервис нижнего уровня;
    \item взаимодействует с одноимённым уровнем другой системы по протоколу.
\end{itemize}

\subsection{Модель OSI}

\begin{center}
\begin{tabular}{cll}
\toprule
Номер & Уровень & Функции \\
\midrule
7 & Прикладной & Сетевые приложения и сервисы \\
6 & Представления & Кодирование, шифрование, сжатие \\
5 & Сеансовый & Управление сеансами связи \\
4 & Транспортный & Доставка между процессами \\
3 & Сетевой & Маршрутизация пакетов \\
2 & Канальный & Передача кадров в пределах канала \\
1 & Физический & Передача битов по среде \\
\bottomrule
\end{tabular}
\end{center}

\subsection{Стек TCP/IP}

Практически в Интернете используется стек TCP/IP:

\begin{center}
\begin{tabular}{ll}
\toprule
Уровень TCP/IP & Примеры протоколов \\
\midrule
Прикладной & HTTP, SMTP, DNS, FTP, SSH \\
Транспортный & TCP, UDP \\
Сетевой & IP, ICMP \\
Канальный & Ethernet, Wi-Fi, PPP \\
Физический & Кабель, оптика, радио \\
\bottomrule
\end{tabular}
\end{center}

\subsection{Инкапсуляция}

При передаче данных каждый уровень добавляет свой заголовок:
\[
\text{данные}
\rightarrow
\text{TCP-сегмент}
\rightarrow
\text{IP-пакет}
\rightarrow
\text{Ethernet-кадр}
\rightarrow
\text{биты}.
\]

При получении выполняется обратный процесс --- декапсуляция.

\section{Локальные и глобальные ИВС}

\subsection{Локальные сети}

\begin{definition}
\textbf{Локальная вычислительная сеть} --- сеть, охватывающая ограниченную территорию: комнату, здание, кампус, предприятие.
\end{definition}

Особенности LAN:
\begin{itemize}
    \item высокая скорость передачи;
    \item малая задержка;
    \item обычно частное владение;
    \item использование Ethernet и Wi-Fi;
    \item ограниченный географический масштаб.
\end{itemize}

Типичные технологии:
\begin{itemize}
    \item Ethernet;
    \item Fast Ethernet;
    \item Gigabit Ethernet;
    \item 10 Gigabit Ethernet;
    \item Wi-Fi.
\end{itemize}

\subsection{Глобальные сети}

\begin{definition}
\textbf{Глобальная вычислительная сеть} --- сеть, объединяющая узлы на больших расстояниях: города, страны, континенты.
\end{definition}

Особенности WAN:
\begin{itemize}
    \item большая территория;
    \item более высокие задержки;
    \item использование операторской инфраструктуры;
    \item маршрутизация через множество промежуточных узлов;
    \item высокая роль протоколов маршрутизации.
\end{itemize}

Примеры WAN:
\begin{itemize}
    \item Интернет;
    \item корпоративные VPN;
    \item магистральные сети операторов;
    \item MPLS-сети.
\end{itemize}

\subsection{Сравнение LAN и WAN}

\begin{center}
\begin{tabular}{lll}
\toprule
Критерий & LAN & WAN \\
\midrule
Масштаб & Здание, кампус & Города, страны, континенты \\
Скорость & Обычно выше & Зависит от оператора \\
Задержка & Низкая & Выше \\
Владение & Частная организация & Операторы связи \\
Технологии & Ethernet, Wi-Fi & IP, MPLS, VPN, оптика \\
\bottomrule
\end{tabular}
\end{center}

\section{Топологии сетей}

\subsection{Физическая и логическая топология}

\begin{definition}
\textbf{Физическая топология} --- реальная схема соединения устройств кабелями или радиоканалами.
\end{definition}

\begin{definition}
\textbf{Логическая топология} --- схема прохождения данных между узлами.
\end{definition}

\subsection{Основные топологии}

\subsubsection{Шина}

Все узлы подключены к общей среде передачи.

Преимущества:
\begin{itemize}
    \item простота;
    \item малые затраты кабеля.
\end{itemize}

Недостатки:
\begin{itemize}
    \item коллизии;
    \item трудность диагностики;
    \item отказ магистрали нарушает работу всей сети.
\end{itemize}

\subsubsection{Звезда}

Все узлы подключены к центральному устройству --- коммутатору.

Преимущества:
\begin{itemize}
    \item простота управления;
    \item отказ одного кабеля не нарушает работу остальных;
    \item хорошая масштабируемость.
\end{itemize}

Недостаток:
\begin{itemize}
    \item отказ центрального коммутатора влияет на всю сеть.
\end{itemize}

\subsubsection{Кольцо}

Каждый узел соединён с двумя соседними.

Преимущества:
\begin{itemize}
    \item упорядоченный доступ к среде;
    \item предсказуемая передача.
\end{itemize}

Недостатки:
\begin{itemize}
    \item сложность перестройки;
    \item чувствительность к отказам без резервирования.
\end{itemize}

\subsubsection{Ячеистая топология}

Узлы имеют множество соединений друг с другом.

Преимущества:
\begin{itemize}
    \item высокая отказоустойчивость;
    \item несколько альтернативных маршрутов.
\end{itemize}

Недостатки:
\begin{itemize}
    \item высокая стоимость;
    \item сложность управления.
\end{itemize}

\section{Технические средства ИВС}

\subsection{Повторитель}

\begin{definition}
\textbf{Повторитель} --- устройство физического уровня, восстанавливающее и усиливающее сигнал.
\end{definition}

Повторитель не анализирует кадры или пакеты. Он работает с битами.

\subsection{Концентратор}

\begin{definition}
\textbf{Концентратор} --- многопортовый повторитель, передающий сигнал со входного порта на все остальные порты.
\end{definition}

Недостаток концентратора --- все устройства находятся в одной области коллизий.

\subsection{Мост}

\begin{definition}
\textbf{Мост} --- устройство канального уровня, соединяющее сегменты сети и фильтрующее кадры по MAC-адресам.
\end{definition}

\subsection{Коммутатор}

\begin{definition}
\textbf{Коммутатор} --- устройство канального уровня, пересылающее кадры между портами на основе таблицы MAC-адресов.
\end{definition}

Принцип работы:
\begin{enumerate}
    \item коммутатор получает кадр;
    \item считывает MAC-адрес источника;
    \item заносит адрес источника в таблицу;
    \item ищет MAC-адрес назначения;
    \item если адрес найден, отправляет кадр на нужный порт;
    \item если адрес неизвестен, рассылает кадр на все порты, кроме входного.
\end{enumerate}

\subsection{Маршрутизатор}

\begin{definition}
\textbf{Маршрутизатор} --- устройство сетевого уровня, пересылающее IP-пакеты между различными сетями на основе таблицы маршрутизации.
\end{definition}

Функции:
\begin{itemize}
    \item выбор маршрута;
    \item разделение широковещательных доменов;
    \item NAT;
    \item фильтрация трафика;
    \item поддержка VPN;
    \item межсетевое взаимодействие.
\end{itemize}

\subsection{Шлюз}

\begin{definition}
\textbf{Шлюз} --- устройство или программа, обеспечивающая взаимодействие сетей или приложений с различными протоколами.
\end{definition}

Пример: почтовый шлюз между разными системами электронной почты.

\subsection{Точка доступа}

\begin{definition}
\textbf{Точка доступа} --- устройство, обеспечивающее подключение беспроводных клиентов к проводной или беспроводной сети.
\end{definition}

\section{Программные средства объединения сетей}

\subsection{Сетевые операционные системы}

Сетевые ОС обеспечивают:
\begin{itemize}
    \item управление пользователями;
    \item сетевую файловую систему;
    \item доступ к принтерам;
    \item удалённое администрирование;
    \item сетевую безопасность;
    \item поддержку протоколов.
\end{itemize}

\subsection{DNS}

\begin{definition}
\textbf{DNS} --- распределённая система доменных имён, преобразующая символьные имена в IP-адреса и наоборот.
\end{definition}

Пример:
\[
\texttt{example.org} \rightarrow 93.184.216.34.
\]

\subsection{DHCP}

\begin{definition}
\textbf{DHCP} --- протокол автоматической выдачи сетевых параметров узлам.
\end{definition}

DHCP выдаёт:
\begin{itemize}
    \item IP-адрес;
    \item маску сети;
    \item адрес шлюза;
    \item DNS-серверы;
    \item срок аренды адреса.
\end{itemize}

\subsection{NAT}

\begin{definition}
\textbf{NAT} --- network address translation, механизм преобразования сетевых адресов.
\end{definition}

Типичный случай:
\[
\text{частный адрес } 192.168.1.10
\rightarrow
\text{публичный адрес маршрутизатора}.
\]

NAT позволяет нескольким внутренним узлам использовать один внешний IP-адрес.

\subsection{VPN}

\begin{definition}
\textbf{VPN} --- virtual private network, виртуальная частная сеть, создающая защищённый логический канал поверх другой сети.
\end{definition}

Назначение VPN:
\begin{itemize}
    \item безопасный удалённый доступ;
    \item объединение филиалов;
    \item шифрование трафика;
    \item туннелирование пакетов.
\end{itemize}

\subsection{VLAN}

\begin{definition}
\textbf{VLAN} --- virtual local area network, виртуальная локальная сеть, позволяющая логически разделить одну физическую сеть на несколько независимых широковещательных доменов.
\end{definition}

Пример:
\[
\text{VLAN 10} = \text{бухгалтерия},\quad
\text{VLAN 20} = \text{разработка},\quad
\text{VLAN 30} = \text{администрация}.
\]

\section{Маршрутизация}

\subsection{Понятие маршрутизации}

\begin{definition}
\textbf{Маршрутизация} --- процесс выбора пути передачи пакетов от источника к получателю через сеть.
\end{definition}

Таблица маршрутизации содержит:
\begin{itemize}
    \item сеть назначения;
    \item маску;
    \item следующий узел;
    \item интерфейс;
    \item метрику.
\end{itemize}

\subsection{Статическая и динамическая маршрутизация}

\textbf{Статическая маршрутизация} задаётся администратором вручную.

Преимущества:
\begin{itemize}
    \item простота;
    \item предсказуемость;
    \item отсутствие служебного трафика.
\end{itemize}

Недостатки:
\begin{itemize}
    \item плохо масштабируется;
    \item требует ручного изменения при сбоях.
\end{itemize}

\textbf{Динамическая маршрутизация} использует протоколы обмена маршрутной информацией.

Примеры:
\begin{itemize}
    \item RIP;
    \item OSPF;
    \item IS-IS;
    \item BGP.
\end{itemize}

\subsection{Алгоритм Дейкстры для маршрутизации по состоянию каналов}

\subsubsection{Постановка}

Дана сеть в виде графа:
\[
G=(V,E),
\]
где $V$ --- маршрутизаторы, $E$ --- каналы связи, каждому ребру сопоставлена неотрицательная стоимость $w(u,v)$.

Нужно найти кратчайшие пути от исходного маршрутизатора $s$ до всех остальных.

\subsubsection{Шаги алгоритма}

\begin{enumerate}
    \item Для всех вершин задать расстояние:
    \[
    d(v):=\infty.
    \]
    Для источника:
    \[
    d(s):=0.
    \]

    \item Множество посещённых вершин:
    \[
    S:=\varnothing.
    \]

    \item Пока есть непосещённые вершины:
    \begin{enumerate}
        \item выбрать непосещённую вершину $u$ с минимальным $d(u)$;
        \item добавить $u$ в $S$;
        \item для каждого соседа $v$ вершины $u$ выполнить релаксацию:
        \[
        \text{если } d(v)>d(u)+w(u,v),
        \text{ то } d(v):=d(u)+w(u,v).
        \]
    \end{enumerate}
\end{enumerate}

\subsubsection{Минимальный пример}

Пусть есть граф:
\[
A \xleftrightarrow{1} B,\quad
A \xleftrightarrow{4} C,\quad
B \xleftrightarrow{2} C,\quad
B \xleftrightarrow{5} D,\quad
C \xleftrightarrow{1} D.
\]

Найдём кратчайшие пути из $A$.

Инициализация:
\[
d(A)=0,\quad d(B)=\infty,\quad d(C)=\infty,\quad d(D)=\infty.
\]

Шаг 1: выбираем $A$.
\[
d(B)=1,\quad d(C)=4.
\]

Шаг 2: выбираем $B$, так как $d(B)=1$.
\[
d(C)=\min(4,1+2)=3,
\]
\[
d(D)=\min(\infty,1+5)=6.
\]

Шаг 3: выбираем $C$, так как $d(C)=3$.
\[
d(D)=\min(6,3+1)=4.
\]

Шаг 4: выбираем $D$.

Итог:
\[
d(A)=0,\quad d(B)=1,\quad d(C)=3,\quad d(D)=4.
\]

Кратчайший путь из $A$ в $D$:
\[
A \rightarrow B \rightarrow C \rightarrow D.
\]

\subsection{Дистанционно-векторная маршрутизация}

Дистанционно-векторные протоколы основаны на обмене таблицами расстояний между соседями.

Для узла $x$ расстояние до назначения $y$ вычисляется как:
\[
D_x(y)=\min_{v \in N(x)} \{c(x,v)+D_v(y)\},
\]
где $N(x)$ --- множество соседей узла $x$, $c(x,v)$ --- стоимость канала между $x$ и $v$.

\subsubsection{Шаги алгоритма}

\begin{enumerate}
    \item Каждый маршрутизатор знает стоимость каналов до своих соседей.
    \item Для себя расстояние равно нулю.
    \item До неизвестных сетей расстояние считается бесконечным.
    \item Периодически каждый маршрутизатор рассылает соседям свой вектор расстояний.
    \item Получив вектор соседа, маршрутизатор пересчитывает свои расстояния по формуле Беллмана--Форда.
    \item Если таблица изменилась, маршрутизатор распространяет обновление дальше.
\end{enumerate}

\section{Коммутация}

\subsection{Коммутация каналов}

При коммутации каналов перед обменом данными устанавливается выделенный путь.

Пример: классическая телефонная сеть.

Преимущества:
\begin{itemize}
    \item гарантированная полоса;
    \item стабильная задержка.
\end{itemize}

Недостатки:
\begin{itemize}
    \item неэффективность при паузах;
    \item необходимость предварительного соединения.
\end{itemize}

\subsection{Коммутация пакетов}

При коммутации пакетов данные разбиваются на пакеты, каждый из которых передаётся независимо или в рамках логического соединения.

Преимущества:
\begin{itemize}
    \item эффективное использование канала;
    \item гибкая маршрутизация;
    \item устойчивость к отказам.
\end{itemize}

Недостатки:
\begin{itemize}
    \item переменная задержка;
    \item возможные потери;
    \item необходимость буферизации.
\end{itemize}

\section{Надёжность и безопасность ИВС}

\subsection{Основные угрозы}

\begin{itemize}
    \item перехват трафика;
    \item подмена адресов;
    \item отказ в обслуживании;
    \item несанкционированный доступ;
    \item вредоносное ПО;
    \item атаки на маршрутизацию;
    \item утечки данных.
\end{itemize}

\subsection{Средства защиты}

\begin{itemize}
    \item межсетевые экраны;
    \item системы обнаружения и предотвращения вторжений;
    \item шифрование;
    \item аутентификация;
    \item сегментация сети;
    \item VPN;
    \item списки контроля доступа;
    \item журналирование;
    \item резервирование.
\end{itemize}

\subsection{Отказоустойчивость}

Методы повышения отказоустойчивости:
\begin{itemize}
    \item резервные каналы связи;
    \item резервные маршрутизаторы;
    \item кластеризация серверов;
    \item репликация данных;
    \item резервное копирование;
    \item балансировка нагрузки;
    \item мониторинг состояния.
\end{itemize}

\section{Связь сетей и распределённых вычислений}

Информационно-вычислительные сети являются инфраструктурной основой распределённых систем и кластеров.

\[
\text{ИВС}
\rightarrow
\text{связь узлов}
\rightarrow
\text{распределённые сервисы}
\rightarrow
\text{кластеры и облака}.
\]

Сетевые протоколы обеспечивают транспорт данных, а распределённые алгоритмы обеспечивают:
\begin{itemize}
    \item согласование состояния;
    \item выбор координатора;
    \item репликацию;
    \item обнаружение отказов;
    \item балансировку нагрузки;
    \item синхронизацию событий.
\end{itemize}

\section{Краткое итоговое сопоставление}

\begin{center}
\begin{tabular}{p{4cm}p{10cm}}
\toprule
Понятие & Смысл \\
\midrule
Многопроцессорная система & Несколько процессоров в одной системе, часто с общей памятью \\
Многомашинный комплекс & Несколько самостоятельных машин, связанных сетью \\
Кластер & Группа узлов, работающих как единый вычислительный ресурс \\
Распределённая система & Независимые компьютеры, совместно предоставляющие сервис \\
ИВС & Сеть технических и программных средств для обработки и передачи информации \\
LAN & Локальная сеть ограниченного масштаба \\
WAN & Глобальная сеть большого масштаба \\
Маршрутизатор & Устройство сетевого уровня для связи разных сетей \\
Коммутатор & Устройство канального уровня для пересылки кадров внутри сети \\
VPN & Защищённая виртуальная сеть поверх другой сети \\
\bottomrule
\end{tabular}
\end{center}

\end{document}
